\ifdefined\inMainDoc\else
\documentclass[12pt,a4paper]{article}
% ================================================================
%  PRÉAMBULE COMMUN - ANNALES CORRIGÉES
%  Université de Kara - FaST - Licence Fondamentale MPC
%  Design optimisé impression noir et blanc
% ================================================================

% -- ENCODAGE & LANGUE --
\usepackage[utf8]{inputenc}
\usepackage[T1]{fontenc}
\usepackage[french]{babel}
\usepackage{lmodern}
\usepackage{microtype}

% -- MISE EN PAGE --
\usepackage[a4paper, top=2.8cm, bottom=2.5cm, left=2.5cm, right=2.5cm, headheight=14pt]{geometry}
\usepackage{setspace}
\setstretch{1.2}
\usepackage{parskip}
\setlength{\parindent}{1.2em}
\setlength{\parskip}{4pt}

% -- MATHS --
\usepackage{amsmath, amssymb, mathtools, amsthm}
\usepackage{mathrsfs}
\usepackage{dsfont}
\usepackage{bm}
\usepackage{cancel}
\usepackage{textcomp}
% Définition de \oiint (intégrale de surface fermée) sans package supplémentaire
\newcommand{\oiint}{\mathop{\iint\!\!\!\!\!\!\!\bigcirc}\limits}

% -- PHYSIQUE & CHIMIE --
\usepackage{physics}
\usepackage{siunitx}
% Lorsque physics est chargé, siunitx omet \qty. On le restaure ici :
\AtBeginDocument{\RenewCommandCopy\qty\SI}
\sisetup{locale = FR, detect-all, output-decimal-marker={,}}
% Commande de substitution pour les unités posant problème avec pdflatex
% (ohm, degreeCelsius, micro, angstrom, percent utilisent des caractères non-ASCII)
\newcommand{\qtyohm}[1]{\ensuremath{\num{#1}\,\Omega}}
\newcommand{\qtydegc}[1]{\ensuremath{\num{#1}\,{}^\circ\text{C}}}
\newcommand{\qtyangstrom}[1]{\ensuremath{\num{#1}\,\text{\AA}}}
\newcommand{\qtypercent}[1]{\ensuremath{\num{#1}\,\%}}
\newcommand{\qtyum}[1]{\ensuremath{\num{#1}\,\mu\text{m}}}
\usepackage[version=4]{mhchem}

% -- GRAPHISMES --
\usepackage{graphicx}
\usepackage{float}
\usepackage{caption}
\usepackage{subcaption}
\usepackage{wrapfig}
\usepackage{tikz}
\usetikzlibrary{calc, arrows.meta, patterns, shapes.geometric}
\usepackage{pgfplots}
\pgfplotsset{compat=1.18}

% -- TABLEAUX --
\usepackage{booktabs}
\usepackage{array}
\usepackage{tabularx}
\usepackage{multirow}

% -- LISTES --
\usepackage{enumitem}
\setlist[enumerate]{topsep=3pt, itemsep=2pt, parsep=0pt}
\setlist[itemize]{topsep=3pt, itemsep=2pt, parsep=0pt, label={.}}

% -- BOÎTES (noir et blanc) --
\usepackage{mdframed}
\usepackage{tcolorbox}
\tcbuselibrary{skins, breakable, theorems}

% Boîte EXERCICE : encadré avec filet épais, fond blanc
\newmdenv[
  linewidth=1.2pt,
  linecolor=black,
  roundcorner=3pt,
  skipabove=10pt,
  skipbelow=6pt,
  innertopmargin=8pt,
  innerbottommargin=8pt,
  innerleftmargin=10pt,
  innerrightmargin=10pt,
  backgroundcolor=white,
  frametitle={\bfseries\small Exercice},
  frametitlealignment=\raggedright,
  frametitleaboveskip=4pt,
  frametitlebelowskip=4pt,
]{boiteexercice}

% Boîte CORRIGÉ : fond gris très clair + filet fin
\newmdenv[
  linewidth=0.6pt,
  linecolor=black,
  leftline=true,
  rightline=false,
  topline=false,
  bottomline=false,
  linecolor=black,
  innerleftmargin=12pt,
  innerrightmargin=8pt,
  innertopmargin=6pt,
  innerbottommargin=6pt,
  backgroundcolor=gray!8,
  skipabove=4pt,
  skipbelow=8pt,
]{boitecorrige}

% Boîte RÉSUMÉ DE COURS : double encadré
\newmdenv[
  linewidth=0.8pt,
  linecolor=black,
  roundcorner=2pt,
  backgroundcolor=gray!12,
  innertopmargin=10pt,
  innerbottommargin=10pt,
  innerleftmargin=12pt,
  innerrightmargin=12pt,
  skipabove=12pt,
  skipbelow=8pt,
]{boiteresume}

% Boîte ATTENTION/REMARQUE : barre gauche épaisse
\newmdenv[
  leftline=true,
  rightline=false,
  topline=false,
  bottomline=false,
  linewidth=3pt,
  linecolor=black,
  backgroundcolor=gray!5,
  innerleftmargin=12pt,
  innerrightmargin=8pt,
  innertopmargin=5pt,
  innerbottommargin=5pt,
  skipabove=6pt,
  skipbelow=6pt,
]{boiteremarque}

% -- DIFFICULTÉ DES EXERCICES --
\newcommand{\facile}{$\star$}
\newcommand{\moyen}{$\star\!\star$}
\newcommand{\difficile}{$\star\!\star\!\star$}
\newcommand{\niveau}[1]{\hfill\textbf{#1}\hspace{2pt}}

% -- EN-TÊTES ET PIEDS DE PAGE --
\usepackage{fancyhdr}
\pagestyle{fancy}
\fancyhf{}
\renewcommand{\headrulewidth}{0.5pt}
\renewcommand{\footrulewidth}{0.3pt}

% -- HYPERLIENS --
\usepackage[hidelinks]{hyperref}
\hypersetup{
    colorlinks=false,
    pdftitle={Annale Corrigée - Licence MPC - Université de Kara}
}

% -- TITRES --
\usepackage{titlesec}
\titleformat{\section}{\Large\bfseries}{\thesection.}{0.8em}{}[\vspace{2pt}\hrule\vspace{4pt}]
\titleformat{\subsection}{\large\bfseries}{\thesubsection.}{0.7em}{}
\titleformat{\subsubsection}{\normalsize\bfseries}{\thesubsubsection.}{0.6em}{}

% -- THÉORÈMES --
\theoremstyle{definition}
\newtheorem{defn}{Définition}[section]
\newtheorem{prop}{Propriété}[section]
\newtheorem{thm}{Théorème}[section]
\newtheorem{corol}{Corollaire}[section]
\newtheorem{exemple}{Exemple}[section]

% -- COMMANDES MATHÉMATIQUES UTILES --
\newcommand{\vect}[1]{\overrightarrow{#1}}
\newcommand{\R}{\mathbb{R}}
\newcommand{\N}{\mathbb{N}}
\newcommand{\Z}{\mathbb{Z}}
\newcommand{\C}{\mathbb{C}}
\renewcommand{\d}{\mathrm{d}}
\newcommand{\deriv}[2]{\dfrac{\d #1}{\d #2}}
\newcommand{\dpart}[2]{\dfrac{\partial #1}{\partial #2}}
\newcommand{\rot}{\overrightarrow{\mathrm{rot}}\,}
\renewcommand{\div}{\mathrm{div}\,}
\renewcommand{\grad}{\overrightarrow{\mathrm{grad}}\,}
\newcommand{\e}{\mathrm{e}}
\newcommand{\ii}{\mathrm{i}}
\renewcommand{\Tr}{\mathrm{Tr}}

% Numérotation des équations par section
\numberwithin{equation}{section}

% -- RÉSULTATS ENCADRÉS --
\newcommand{\res}[1]{\[\boxed{#1}\]}

% -- REMARQUE INLINE --
\newcommand{\rmq}[1]{\begin{boiteremarque}\textbf{Remarque.} #1\end{boiteremarque}}

% -- MISE EN VALEUR DU RÉSULTAT FINAL --
\newcommand{\reponse}[1]{%
  \begin{center}
    \fbox{\begin{minipage}{0.92\textwidth}\centering #1\end{minipage}}
  \end{center}%
}


\lhead{\small\textit{Physique Mathématique - Licence 3}}
\rhead{\small\textit{FaST, Université de Kara}}
\cfoot{\thepage}

\begin{document}
\fi

\begin{center}
  {\LARGE\bfseries PHYSIQUE MATHÉMATIQUE}\\[0.3cm]
  {\normalsize Licence 3 - Physique}\\[0.1cm]
  \rule{0.7\textwidth}{0.8pt}
\end{center}

\section*{Résumé du cours}

\begin{boiteresume}

\textbf{1. Formes différentielles de degré 1.} Une forme différentielle $\omega = P(x,y)\,\d x + Q(x,y)\,\d y$ est \textit{fermée} si $\frac{\partial P}{\partial y} = \frac{\partial Q}{\partial x}$, et \textit{exacte} s'il existe $f$ (une primitive) telle que $\d f = \omega$. Toute forme exacte est fermée. Sur un ouvert étoilé, la réciproque est vraie (théorème de Poincaré). L'intégrale curviligne d'une forme exacte le long d'un chemin ne dépend que des extrémités.

\textbf{2. Formule de Green-Riemann.} Pour une forme $\omega = P\,\d x + Q\,\d y$ de classe $\mathcal{C}^1$ sur un domaine $K$ de frontière $\gamma$ orientée positivement :
\[
\oint_\gamma P\,\d x + Q\,\d y = \iint_K \left(\frac{\partial Q}{\partial x} - \frac{\partial P}{\partial y}\right)\d x\,\d y
\]
Application : aire $A = \frac{1}{2}\oint_\gamma x\,\d y - y\,\d x$.

\textbf{3. Équations aux dérivées partielles (EDP) du 2\textsuperscript{e} ordre.} Pour $au_{xx}+2bu_{xy}+cu_{yy}+\cdots=0$, le discriminant $\Delta = b^2 - ac$ détermine le type : $\Delta>0$ (hyperbolique), $\Delta=0$ (parabolique), $\Delta<0$ (elliptique). La réduction à la forme canonique se fait par un changement de variables basé sur les courbes caractéristiques.

\textbf{4. Transformée de Fourier.} $\hat{f}(\xi) = \int_{-\infty}^{+\infty} f(x)\,\e^{-2\pi\ii x\xi}\,\d x$. Propriétés : linéarité, décalage, dérivation ($\widehat{f'} = 2\pi\ii\xi\,\hat{f}$), convolution ($\widehat{f*g} = \hat{f}\cdot\hat{g}$).

\textbf{5. Intégrale de Fourier.} Théorème de Parseval : $\int_{-\infty}^{+\infty}|f(x)|^2\,\d x = \int_{-\infty}^{+\infty}|\hat{f}(\xi)|^2\,\d\xi$.

\end{boiteresume}

\newpage
\section{Formes différentielles}

\subsection*{Exercice 1 - Forme exacte sur un ouvert étoilé \niveau{\facile}}

\begin{boiteexercice}
Soit $\omega = \dfrac{x\,\d y - y\,\d x}{x^2+y^2}$ définie sur $U=\{(x,y)\in\mathbb{R}^2,\;x>0\}$.
\begin{enumerate}
  \item Montrer que $\omega$ est fermée sur $U$.
  \item Montrer que $U$ est étoilé, puis en déduire que $\omega$ est exacte.
  \item Déterminer toutes les primitives de $\omega$ sur $U$.
\end{enumerate}
\end{boiteexercice}

\begin{boitecorrige}
\textbf{1.} Posons $P = \frac{-y}{x^2+y^2}$ et $Q = \frac{x}{x^2+y^2}$. On calcule :
\[
\frac{\partial P}{\partial y} = \frac{-(x^2+y^2)+2y^2}{(x^2+y^2)^2} = \frac{y^2-x^2}{(x^2+y^2)^2}
\]
\[
\frac{\partial Q}{\partial x} = \frac{(x^2+y^2)-2x^2}{(x^2+y^2)^2} = \frac{y^2-x^2}{(x^2+y^2)^2}
\]
Les dérivées croisées sont égales : $\omega$ est fermée sur $U$.

\textbf{2.} $U$ est étoilé (convexe, donc étoilé par rapport à tout point, par exemple $(1,0)$). Par le théorème de Poincaré, toute forme fermée sur un ouvert étoilé est exacte. Donc $\omega$ est exacte sur $U$.

\textbf{3.} On cherche $f$ telle que $\partial f/\partial x = -y/(x^2+y^2)$ et $\partial f/\partial y = x/(x^2+y^2)$.

En intégrant la seconde équation par rapport à $y$ :
\[
f(x,y) = \arctan\!\left(\frac{y}{x}\right) + H(x)
\]
On substitue dans la première : $\frac{\partial}{\partial x}\arctan(y/x) = \frac{-y/x^2}{1+(y/x)^2} = \frac{-y}{x^2+y^2}$, donc $H'(x)=0$, $H=\text{cste}$.

Les primitives sont $f(x,y) = \arctan(y/x) + C$, $C\in\mathbb{R}$.
\end{boitecorrige}

\subsection*{Exercice 2 - Primitive d'une forme exacte \niveau{\moyen}}

\begin{boiteexercice}
Soit $\omega = \dfrac{2x}{y}\,\d x - \dfrac{x^2}{y^2}\,\d y$ sur $U=\{(x,y)\in\mathbb{R}^2,\;y>0\}$.
\begin{enumerate}
  \item Montrer que $\omega$ est exacte sur $U$.
  \item Déterminer une primitive de $\omega$.
  \item Calculer $\int_C \omega$ où $C$ est un chemin de classe $\mathcal{C}^1$ allant de $A(1,2)$ à $B(3,8)$.
\end{enumerate}
\end{boiteexercice}

\begin{boitecorrige}
\textbf{1.} $\frac{\partial P}{\partial y} = \frac{2x \cdot (-1)}{y^2} = \frac{-2x}{y^2}$ et $\frac{\partial Q}{\partial x} = \frac{-2x}{y^2}$. Formes fermées. $U$ est étoilé (convexe), donc $\omega$ est exacte.

\textbf{2.} On intègre $\partial f/\partial x = 2x/y$ : $f(x,y) = x^2/y + H(y)$. On vérifie $\partial f/\partial y = -x^2/y^2 + H'(y) = -x^2/y^2$, donc $H'(y)=0$, $H=\text{cste}$.

Primitive : $f(x,y) = \dfrac{x^2}{y}$.

\textbf{3.} Comme $\omega$ est exacte et $f$ en est une primitive :
\[
\int_C \omega = f(B) - f(A) = \frac{9}{8} - \frac{1}{2} = \frac{5}{8}
\]
\end{boitecorrige}

\subsection*{Exercice 3 - Facteur intégrant \niveau{\moyen}}

\begin{boiteexercice}
Soit $\omega = (x^2+y^2-a^2)\,\d x - 2ay\,\d y$ sur $\mathbb{R}^2$, où $a\neq0$.
\begin{enumerate}
  \item Montrer que $\omega$ n'est pas exacte.
  \item Chercher une fonction $f:\mathbb{R}\to\mathbb{R}$ de classe $\mathcal{C}^1$ telle que $\alpha = f(x)\omega$ soit exacte sur $\mathbb{R}^2$.
  \item Calculer une primitive de $\alpha$.
  \item En déduire $\int_C \alpha$ sur le cercle de rayon $R$ centré à l'origine.
\end{enumerate}
\end{boiteexercice}

\begin{boitecorrige}
\textbf{1.} $P = x^2+y^2-a^2$, $Q=-2ay$. $\partial P/\partial y = 2y \neq \partial Q/\partial x = 0$. Non exacte.

\textbf{2.} On pose $P_1 = f(x)(x^2+y^2-a^2)$ et $Q_1 = -2ayf(x)$. Pour que $\alpha$ soit fermée :
\[
\frac{\partial P_1}{\partial y} = 2yf(x) = \frac{\partial Q_1}{\partial x} = -2ayf'(x)
\]
Donc $f'(x) = -\frac{f(x)}{a}$, d'où $f(x) = \e^{-x/a}$. La forme est fermée sur l'ouvert étoilé $\mathbb{R}^2$, donc exacte.

\textbf{3.} On cherche $F$ telle que $\partial F/\partial y = -2ay\e^{-x/a}$ : $F = -ay^2\e^{-x/a} + H(x)$. En substituant : $H'(x) = (x^2-a^2)\e^{-x/a}$. Par intégrations par parties : $H(x) = (-ax^2 - 2a^2x - a^3)\e^{-x/a}$.

Primitive : $F(x,y) = -(ay^2 + ax^2 + 2a^2x + a^3)\e^{-x/a}$.

\textbf{4.} La forme $\alpha$ est exacte et $C$ est fermée, donc $\int_C \alpha = F(\text{retour}) - F(\text{départ}) = 0$.
\end{boitecorrige}

\subsection*{Exercice 4 - Formule de Green-Riemann \niveau{\moyen}}

\begin{boiteexercice}
Soit $K = \{(x,y)\in\mathbb{R}^2,\;x\geq0,\;y\geq0,\;x^2+y^2\leq1\}$ et $\omega = xy^2\,\d x + 2xy\,\d y$.

Calculer $\int_\gamma \omega$ (où $\gamma = \partial K$ orienté positivement) de deux façons : par une paramétrisation de $\gamma$, puis par la formule de Green-Riemann.
\end{boiteexercice}

\begin{boitecorrige}
\textbf{Méthode 1 - Paramétrisation.} Le bord $\gamma$ se décompose en trois arcs :
\begin{itemize}[nosep]
  \item $C_1$ : $(t,0)$, $t\in[0,1]$ : $\omega|_{C_1}=0$ (car $y=0$).
  \item $C_3$ : $(0,t)$, $t\in[1,0]$ : $\omega|_{C_3}=0$ (car $x=0$).
  \item $C_2$ : $(\cos t,\sin t)$, $t\in[0,\pi/2]$.
\end{itemize}
Sur $C_2$ :
\begin{align*}
\int_{C_2}\omega &= \int_0^{\pi/2}\bigl[\cos t\sin^2 t(-\sin t) + 2\cos t\sin t\cos t\bigr]\d t\\
&= \int_0^{\pi/2}\bigl(-\cos t\sin^3 t + 2\cos^2 t\sin t\bigr)\d t
\end{align*}
\[
= \left[-\frac{\sin^4 t}{4}\right]_0^{\pi/2} + \left[-\frac{2\cos^3 t}{3}\right]_0^{\pi/2}
= -\frac{1}{4} + \frac{2}{3} = \frac{5}{12}
\]

\textbf{Méthode 2 - Green-Riemann.}
\[
\frac{\partial Q}{\partial x} - \frac{\partial P}{\partial y} = 2y - 2xy
\]
En coordonnées polaires ($x=r\cos\theta$, $y=r\sin\theta$, $r\in[0,1]$, $\theta\in[0,\pi/2]$) :
\begin{align*}
\iint_K(2y-2xy)\,\d x\,\d y &= \int_0^1\!\int_0^{\pi/2}(2r\sin\theta - 2r^2\sin\theta\cos\theta)\,r\,\d\theta\,\d r\\
&= \int_0^1(2r^2-r^3)\,\d r = \frac{2}{3}-\frac{1}{4} = \frac{5}{12}
\end{align*}
Les deux méthodes donnent bien $\dfrac{5}{12}$.
\end{boitecorrige}

\section{Équations aux dérivées partielles}

\subsection*{Exercice 5 - EDP parabolique \niveau{\difficile}}

\begin{boiteexercice}
Considérer le problème de Cauchy :
\[
\begin{cases}
\dfrac{\partial^2 u}{\partial x^2} - 2\dfrac{\partial^2 u}{\partial x\partial y} + \dfrac{\partial^2 u}{\partial y^2} + \dfrac{\partial u}{\partial x} - \dfrac{\partial u}{\partial y} = 0\\[4pt]
u(x,0)=f(x),\quad \dfrac{\partial u}{\partial y}(x,0)=g(x)
\end{cases}
\]
\begin{enumerate}
  \item Indiquer le type de l'EDP.
  \item Réduire à la forme canonique avec $\eta=y$.
  \item Déterminer la solution générale.
  \item Résoudre le PCIs. Traiter les cas $f=g=1$ et $f(x)=x^2$, $g(x)=0$.
\end{enumerate}
\end{boiteexercice}

\begin{boitecorrige}
\textbf{1.} Ici $a=1$, $b=-1$, $c=1$, donc $\Delta = b^2-ac = 1-1=0$. EDP \textbf{parabolique}.

L'équation caractéristique est $(\d y + \d x)^2 = 0$, soit $x+y = \text{cste}$.

\textbf{2.} On pose $\xi = x+y$ et $\eta = y$. Les dérivées deviennent :
\[
\frac{\partial u}{\partial x} = \frac{\partial v}{\partial\xi},\quad \frac{\partial u}{\partial y} = \frac{\partial v}{\partial\xi}+\frac{\partial v}{\partial\eta},\quad \frac{\partial^2 u}{\partial y^2} = \frac{\partial^2 v}{\partial\xi^2}+2\frac{\partial^2 v}{\partial\xi\partial\eta}+\frac{\partial^2 v}{\partial\eta^2}
\]
En substituant dans l'EDP, les termes en $\partial^2v/\partial\xi^2$ et $\partial^2v/\partial\xi\partial\eta$ s'annulent. La forme canonique est :
\[
\frac{\partial^2 v}{\partial\eta^2} - \frac{\partial v}{\partial\eta} = 0
\]

\textbf{3.} En posant $w=\partial v/\partial\eta$, on obtient $\partial w/\partial\eta = w$, d'où $w = \Phi(\xi)\e^\eta$. En intégrant : $v(\xi,\eta)=\Phi(\xi)\e^\eta+\Psi(\xi)$. La solution générale est :
\[
u(x,y) = \Phi(x+y)\e^y + \Psi(x+y)
\]

\textbf{4.} Conditions initiales :
$u(x,0) = \Phi(x)+\Psi(x) = f(x)$.
$\partial u/\partial y(x,0) = \Phi'(x)+\Phi(x)+\Psi'(x) = g(x)$.

Résolution : $[\Phi+\Psi]'=f'(x)$ et $\Phi'+\Phi+\Psi'=g(x)$, soit $\Phi = g-f'$ et $\Psi = f-g+f'$.

Solution du PCIs : $u(x,y) = [g(x+y)-f'(x+y)]\e^y + f(x+y)-g(x+y)+f'(x+y)$.

\textit{Cas 1 :} $f=g=1$, $f'=0$. $u(x,y) = \e^y$.

\textit{Cas 2 :} $f(x)=x^2$, $g(x)=0$, $f'(x)=2x$. $u(x,y) = -2(x+y)\e^y + (x+y)^2 + 2(x+y) = (x+y)[(x+y)+2(1-\e^y)]$.
\end{boitecorrige}

\subsection*{Exercice 6 - Aire par Green-Riemann \niveau{\moyen}}

\begin{boiteexercice}
Calculer l'aire du domaine délimité par un quart de l'astroïde $x=a\cos^3 t$, $y=a\sin^3 t$ pour $t\in[0,\pi/2]$, et les deux axes de coordonnées, en utilisant la formule de Green-Riemann.
\end{boiteexercice}

\begin{boitecorrige}
L'aire s'exprime par $A = \frac{1}{2}\oint_\gamma x\,\d y - y\,\d x$ où $\gamma$ est le contour orienté positivement.

Sur l'arc de l'astroïde ($t$ de 0 à $\pi/2$) : $\d x = -3a\cos^2 t\sin t\,\d t$, $\d y = 3a\sin^2 t\cos t\,\d t$.

\[
A = \frac{1}{2}\int_0^{\pi/2}\left[a\cos^3 t\cdot3a\sin^2 t\cos t - a\sin^3 t\cdot(-3a\sin t\cos^2 t)\right]\d t
\]
\[
= \frac{3a^2}{2}\int_0^{\pi/2}\left(\cos^4 t\sin^2 t + \sin^4 t\cos^2 t\right)\d t
= \frac{3a^2}{2}\int_0^{\pi/2}\cos^2 t\sin^2 t\,\d t
\]
\[
= \frac{3a^2}{2}\int_0^{\pi/2}\frac{\sin^2(2t)}{4}\,\d t = \frac{3a^2}{8}\int_0^{\pi/2}\frac{1-\cos(4t)}{2}\,\d t = \frac{3a^2}{16}\left[\frac{\pi}{2}\right] = \frac{3\pi a^2}{32}
\]

L'aire d'un quart de l'astroïde est $\dfrac{3\pi a^2}{32}$, donc l'aire totale est $A_{total} = \dfrac{3\pi a^2}{8}$.
\end{boitecorrige}


\subsection*{Exercice 7 - Forme différentielle exacte et intégrale curviligne \niveau{\moyen}}

\begin{boiteexercice}
Soit $\omega = (y^3 - 6xy^2)\,\d x + (3xy^2 - 6x^2y)\,\d y$. Montrer que $\omega$ est une forme différentielle exacte sur $\mathbb{R}^2$. En déduire l'intégrale curviligne le long du demi-cercle supérieur de diamètre $[AB]$ de $A(1,2)$ vers $B(3,4)$.
\end{boiteexercice}

\begin{boitecorrige}
On pose $P = y^3 - 6xy^2$ et $Q = 3xy^2 - 6x^2y$.
\[
\frac{\partial P}{\partial y} = 3y^2 - 12xy = \frac{\partial Q}{\partial x}
\]
Donc $\omega$ est fermée. $\mathbb{R}^2$ est étoilé, donc $\omega$ est exacte (théorème de Poincaré).

On cherche une primitive $f$ telle que $\partial f/\partial x = P$ et $\partial f/\partial y = Q$.
En intégrant la première : $f(x,y) = xy^3 - 3x^2y^2 + H(y)$.
En substituant dans la seconde : $3xy^2 - 6x^2y + H'(y) = 3xy^2 - 6x^2y$, donc $H'(y) = 0$, $H = \text{cste}$.

Primitive : $f(x,y) = xy^3 - 3x^2y^2$.

Comme $\omega$ est exacte, l'intégrale ne dépend que des extrémités :
\begin{align*}
\int_C \omega &= f(B) - f(A) = f(3,4) - f(1,2)\\
&= (3\cdot64 - 3\cdot9\cdot16) - (1\cdot8 - 3\cdot4)\\
&= (192 - 432) - (8-12) = -240 + 4 = -236
\end{align*}
\end{boitecorrige}

\subsection*{Exercice 8 - Intégrale le long d'une cardioïde \niveau{\moyen}}

\begin{boiteexercice}
Soit $\omega = (x+y)\,\d x + (x-y)\,\d y$. Calculer l'intégrale curviligne le long de la demi-cardioïde d'équation polaire $r = 1 + \cos\theta$ avec $\theta$ allant de $0$ à $\pi$.
\end{boiteexercice}

\begin{boitecorrige}
On vérifie que $\omega$ est exacte : $\partial P/\partial y = 1 = \partial Q/\partial x$. La forme est fermée sur $\mathbb{R}^2$ étoilé, donc exacte.

On cherche une primitive : en intégrant $\partial f/\partial x = x+y$, on obtient $f(x,y) = \frac{x^2}{2} + xy + H(y)$. Puis $\partial f/\partial y = x + H'(y) = x - y$, donc $H'(y) = -y$, $H(y) = -y^2/2$.

Primitive : $f(x,y) = \dfrac{x^2}{2} + xy - \dfrac{y^2}{2}$.

La demi-cardioïde commence en $\theta=0$ : $r=2$, point $(2,0)$, et finit en $\theta=\pi$ : $r=0$, point $(0,0)$.
\[
\int_C \omega = f(0,0) - f(2,0) = 0 - \frac{4}{2} = -2
\]
\end{boitecorrige}

\newpage
\section{Transformée de Fourier et de Laplace}

\subsection*{Exercice 9 - Transformée de Laplace \niveau{\moyen}}

\begin{boiteexercice}
Calculer les transformées de Laplace suivantes :
\begin{enumerate}[label=(\alph*)]
  \item $\mathcal{L}\!\left[(t^2+t-\e^{-3t})\mathbb{H}(t)\right]$
  \item $\mathcal{L}\!\left[(t+2)\mathbb{H}(t)+(t+3)\mathbb{H}(t-2)\right]$
  \item $\mathcal{L}\!\left[(t^2+t+1)\e^{-2t}\mathbb{H}(t)\right]$
\end{enumerate}
\end{boiteexercice}

\begin{boitecorrige}
\textbf{(a)} Par linéarité :
\[
F(p) = \frac{2}{p^3} + \frac{1}{p^2} - \frac{1}{p+3}
\]

\textbf{(b)} On réécrit $(t+3)\mathbb{H}(t-2) = ((t-2)+5)\mathbb{H}(t-2)$ :
\begin{align*}
\mathcal{L}[(t+2)\mathbb{H}(t)] &= \frac{1}{p^2} + \frac{2}{p} \\
\mathcal{L}[((t-2)+5)\mathbb{H}(t-2)] &= \left(\frac{1}{p^2}+\frac{5}{p}\right)\e^{-2p}
\end{align*}
\[
F(p) = \frac{2p+1}{p^2} + \left(\frac{1+5p}{p^2}\right)\e^{-2p}
\]

\textbf{(c)} Par le théorème du glissement ($\mathcal{L}[f(t)\e^{at}] = F(p-a)$) :
\[
\mathcal{L}[(t^2+t+1)\mathbb{H}(t)] = \frac{2}{p^3}+\frac{1}{p^2}+\frac{1}{p}
\]
Donc $F(p) = \dfrac{(p+2)^2+(p+2)+2}{(p+2)^3} = \dfrac{p^2+5p+8}{(p+2)^3}$.
\end{boitecorrige}

\subsection*{Exercice 10 - Transformée de Laplace inverse \niveau{\moyen}}

\begin{boiteexercice}
Calculer les originaux suivants :
\begin{enumerate}[label=(\alph*)]
  \item $\mathcal{L}^{-1}\!\left[\dfrac{p+2}{(p+3)(p+4)}\right]$
  \item $\mathcal{L}^{-1}\!\left[\dfrac{3}{(p+5)^2}\right]$
  \item $\mathcal{L}^{-1}\!\left[\dfrac{p-1}{p^2+2p+5}\right]$
\end{enumerate}
\end{boiteexercice}

\begin{boitecorrige}
\textbf{(a)} Décomposition en éléments simples : $\dfrac{p+2}{(p+3)(p+4)} = \dfrac{-1}{p+3}+\dfrac{2}{p+4}$. D'où :
\[
f(t) = (-\e^{-3t}+2\e^{-4t})\mathbb{H}(t)
\]

\textbf{(b)} En utilisant $\mathcal{L}^{-1}[1/p^2] = t$ et le théorème du glissement :
\[
f(t) = 3t\e^{-5t}\mathbb{H}(t)
\]

\textbf{(c)} On complète le carré au dénominateur : $(p+1)^2+4$. Alors :
\[
\frac{p-1}{(p+1)^2+4} = \frac{p+1}{(p+1)^2+4} - \frac{2}{(p+1)^2+4}
\]
\[
f(t) = \e^{-t}(\cos 2t - \sin 2t)\mathbb{H}(t)
\]
\end{boitecorrige}

\subsection*{Exercice 11 - Résolution d'une EDO par Laplace \niveau{\difficile}}

\begin{boiteexercice}
On désigne par $F(p)$ et $G(p)$ les transformées de Laplace de $f(t)$ et $g(t)$ :
\[
F(p) = \frac{3p+7}{(p+2)(p+4)^2}, \quad G(p) = \frac{5p+27}{(p+4)^2+9}
\]
\begin{enumerate}
  \item Déterminer l'original $h(t)$ dont l'image est $[F(p)-G(p)]\e^{-2p}$.
  \item En utilisant la transformée de Laplace, résoudre $y'' + 3y' + 2y = 3\e^{-4t}$, $y(0)=2$, $y'(0)=-1$.
\end{enumerate}
\end{boiteexercice}

\begin{boitecorrige}
\textbf{1.} Décomposition de $F(p)$ en éléments simples :
\[
F(p) = \frac{1/4}{p+2} - \frac{1/4}{p+4} + \frac{5/2}{(p+4)^2}
\]
Donc $f(t) = \frac{1}{4}\e^{-2t} - \frac{1}{4}\e^{-4t} + \frac{5}{2}t\e^{-4t}$.

Pour $G(p) = 5\dfrac{p+4}{(p+4)^2+9} + \dfrac{7}{(p+4)^2+9}$, on obtient :
$g(t) = 5\e^{-4t}\cos(3t) + \dfrac{7}{3}\e^{-4t}\sin(3t)$.

Par le théorème du décalage temporel : $h(t) = [f(t-2)-g(t-2)]\mathbb{H}(t-2)$.

\textbf{2.} Application de la transformée de Laplace avec les conditions initiales :
\[
(s^2+3s+2)Y(s) - 2s - 5 = \frac{3}{s+4}
\]
\[
Y(s) = \frac{2s+5}{(s+1)(s+2)} + \frac{3}{(s+1)(s+2)(s+4)} = \frac{4}{s+1} - \frac{5}{2(s+2)} + \frac{1}{2(s+4)}
\]
\[
\boxed{y(t) = 4\e^{-t} - \frac{5}{2}\e^{-2t} + \frac{1}{2}\e^{-4t}}
\]
\end{boitecorrige}

\subsection*{Exercice 12 - Jacobien d'un changement de variables \niveau{\facile}}

\begin{boiteexercice}
\begin{enumerate}
  \item Calculer la matrice jacobienne et le jacobien du changement de variables en coordonnées polaires :
  $x = r\cos\theta$, $y = r\sin\theta$.
  \item Calculer le jacobien du changement de variables $u = x+y$, $v = \dfrac{x}{x+y}$.
\end{enumerate}
\end{boiteexercice}

\begin{boitecorrige}
\textbf{1.} La matrice jacobienne est :
\[
J(r,\theta) = \begin{pmatrix} \cos\theta & -r\sin\theta \\ \sin\theta & r\cos\theta \end{pmatrix}
\]
\[
\det(J) = r\cos^2\theta + r\sin^2\theta = r
\]

\textbf{2.} La matrice jacobienne de $(u,v)$ par rapport à $(x,y)$ :
\[
J = \begin{pmatrix} 1 & 1 \\ \dfrac{y}{(x+y)^2} & \dfrac{-x}{(x+y)^2} \end{pmatrix}
\]
\[
\det(J) = \frac{-x}{(x+y)^2} - \frac{y}{(x+y)^2} = \frac{-(x+y)}{(x+y)^2} = \frac{-1}{x+y} = \frac{-1}{u}
\]
\end{boitecorrige}

\newpage
\subsection*{Exercice 13 - Fonctions Gamma et Bêta \niveau{\difficile}}

\begin{boiteexercice}
\begin{enumerate}
  \item Démontrer la relation :
  \[
  \int_0^{\pi/2} \cos^{2x-1}\theta\,\sin^{2y-1}\theta\,\d\theta = \frac{\Gamma(x)\Gamma(y)}{2\,\Gamma(x+y)}
  \]
  \item Calculer $\beta(1,n)$ et prouver le résultat par récurrence.
  \item Calculer $\beta(1,1)$ et $\beta\!\left(\tfrac{1}{2},\tfrac{1}{2}\right)$.
\end{enumerate}
\end{boiteexercice}

\begin{boitecorrige}
\textbf{1.} En utilisant $\Gamma(x) = 2\int_0^{+\infty}\e^{-t^2}t^{2x-1}\,\d t$ :
\begin{align*}
\Gamma(x)\Gamma(y) &= 4\int_0^{+\infty}\int_0^{+\infty}\e^{-(t^2+u^2)}t^{2x-1}u^{2y-1}\,\d t\,\d u
\end{align*}
On effectue le changement polaire $t = r\cos\theta$, $u = r\sin\theta$, jacobien $r$ :
\begin{align*}
\Gamma(x)\Gamma(y) &= 4\left(\int_0^{+\infty}\e^{-r^2}r^{2(x+y)-1}\,\d r\right)\left(\int_0^{\pi/2}\cos^{2x-1}\theta\,\sin^{2y-1}\theta\,\d\theta\right) \\
&= 2\,\Gamma(x+y)\int_0^{\pi/2}\cos^{2x-1}\theta\,\sin^{2y-1}\theta\,\d\theta
\end{align*}
D'où la relation cherchée.

\textbf{2.} Pour $x=1$ : $\beta(1,n) = \int_0^1(1-t)^{n-1}\,\d t = \dfrac{1}{n}$.

Par récurrence, en utilisant $\beta(x,y+1) = \frac{y}{x+y}\beta(x,y)$ :
$\beta(1,k+1) = \frac{k}{k+1}\cdot\frac{1}{k} = \frac{1}{k+1}$. Résultat prouvé.

\textbf{3.} $\beta(1,1) = 1$.
\[
\beta\!\left(\tfrac{1}{2},\tfrac{1}{2}\right) = \frac{\Gamma(1/2)^2}{\Gamma(1)} = \frac{(\sqrt{\pi})^2}{1} = \pi
\]
\end{boitecorrige}

\newpage
\subsection*{Exercice 14 - EDP : équation de transport \niveau{\moyen}}

\begin{boiteexercice}
\begin{enumerate}
  \item Montrer que la fonction $u(x,y) = f(x+iy) + g(x-iy)$, où $f$ et $g$ sont deux fois dérivables sur $\mathbb{R}$, est harmonique ($\Delta u = 0$).
  \item Trouver les solutions particulières réelles de l'équation de Laplace correspondant aux cas : (a) $f(x) = x^2$, $g = 0$ ; (b) $f(x) = \e^x$, $g = 0$.
  \item Résoudre $u_t + 2u_x = 0$ avec $u(0,x) = \frac{1}{x^2+1}$. Évaluer $u\!\left(\frac{1}{2},x\right)$.
\end{enumerate}
\end{boiteexercice}

\begin{boitecorrige}
\textbf{1.} On calcule :
\[
\frac{\partial^2 u}{\partial x^2} = f''(x+iy) + g''(x-iy), \quad \frac{\partial^2 u}{\partial y^2} = -\bigl(f''(x+iy) + g''(x-iy)\bigr)
\]
Donc $\Delta u = 0$ : $u$ est harmonique.

\textbf{2.} (a) $f(x+iy) = (x+iy)^2 = x^2 - y^2 + 2ixy$. Les parties réelles et imaginaires sont les solutions $x^2 - y^2$ et $xy$.

(b) $f(x+iy) = \e^x(\cos y + i\sin y)$. Solutions : $\e^x\cos y$ et $\e^x\sin y$.

\textbf{3.} C'est une équation de transport. La solution est $u(t,x) = g(x-2t)$ avec la condition initiale :
\[
u(t,x) = \frac{1}{(x-2t)^2+1}
\]
\[
u\!\left(\tfrac{1}{2},x\right) = \frac{1}{(x-1)^2+1}
\]
\end{boitecorrige}

\newpage
\subsection*{Exercice 15 - EDP : type et classification \niveau{\moyen}}

\begin{boiteexercice}
\begin{enumerate}
  \item Dire si l'EDP $\left(\dfrac{\partial u}{\partial x}\right)^2 - u\,\dfrac{\partial u}{\partial y} = 0$ est linéaire ou non linéaire. Justifier.
  \item Résoudre l'EDP :
  \[
  4\frac{\partial u}{\partial t} - 3\frac{\partial u}{\partial x} = 0, \quad u(0,x) = x^3
  \]
\end{enumerate}
\end{boiteexercice}

\begin{boitecorrige}
\textbf{1.} L'EDP est \textbf{non linéaire}. En effet, si $L$ désigne l'opérateur associé, alors $L(\alpha u) \neq \alpha L(u)$ car le terme $(\partial u/\partial x)^2$ est quadratique en $u$. On vérifie : $L(\alpha u) = \alpha^2(\partial u/\partial x)^2 - \alpha u(\partial u/\partial y) \neq \alpha L(u)$ en général.

\textbf{2.} C'est une équation de transport $u_t + au_x = 0$ avec $a = -3/4$ (après division par 4). La solution générale est $u(t,x) = g(x + \frac{3}{4}t)$. La condition initiale donne $g(x) = x^3$, donc :
\[
\boxed{u(t,x) = \left(x + \frac{3}{4}t\right)^3}
\]
\end{boitecorrige}

\subsection*{Exercice 16 - Résolution par Laplace d'un système \niveau{\difficile}}

\begin{boiteexercice}
Utiliser la transformée de Laplace pour résoudre les équations différentielles suivantes :
\begin{enumerate}[label=(\alph*)]
  \item $x'(t) + x(t) = t\mathbb{H}(t) - t\mathbb{H}(t-1)$, avec $x(0) = 0$.
  \item $x''(t) + x'(t) = \mathbb{H}(t)$, avec $x(0) = 0$, $x'(0) = 0$.
\end{enumerate}
\end{boiteexercice}

\begin{boitecorrige}
\textbf{(a)} Application de la transformée de Laplace :
\[
(p+1)X(p) = \frac{1}{p^2} - \left(\frac{1}{p^2}+\frac{1}{p}\right)\e^{-p}
\]
\[
X(p) = \frac{1}{p^2(p+1)} - \frac{\e^{-p}}{p^2} = \left(\frac{1}{p^2} - \frac{1}{p} + \frac{1}{p+1}\right) - \frac{\e^{-p}}{p^2}
\]
\[
x(t) = (t-1+\e^{-t})\mathbb{H}(t) - (t-1)\mathbb{H}(t-1)
\]

\textbf{(b)} Application de la transformée de Laplace :
\[
(p^2+p)X(p) = \frac{1}{p} \implies X(p) = \frac{1}{p^2(p+1)} = \frac{1}{p^2} - \frac{1}{p} + \frac{1}{p+1}
\]
\[
\boxed{x(t) = (t-1+\e^{-t})\mathbb{H}(t)}
\]
\end{boitecorrige}

\subsection*{Exercice 17 - Intégrale de Fourier et forme canonique d'une EDP \niveau{\difficile}}

\begin{boiteexercice}
Considérer l'EDP linéaire du second ordre :
\[
u_{xx} + 4u_{xy} + 4u_{yy} + u_x + 2u_y = 0
\]
\begin{enumerate}
  \item Déterminer le type (elliptique, parabolique, hyperbolique).
  \item Trouver les équations caractéristiques.
  \item Effectuer le changement de variables approprié et écrire la forme canonique.
  \item Donner la solution générale.
\end{enumerate}
\end{boiteexercice}

\begin{boitecorrige}
\textbf{1.} Ici $a = 1$, $b = 2$, $c = 4$. Le discriminant : $\Delta = b^2 - ac = 4 - 4 = 0$. L'EDP est \textbf{parabolique}.

\textbf{2.} L'équation caractéristique est $(dy)^2 - 2\cdot 2\,dx\,dy + 1\cdot(dx)^2 = 0$, soit $(dy - 2\,dx)^2 = 0$, d'où $y - 2x = \text{cste}$.

\textbf{3.} Changement de variables : $\xi = y - 2x$, $\eta = y$. Les dérivées partielles deviennent :
\[
u_x = -2v_\xi, \quad u_y = v_\xi + v_\eta, \quad u_{xx} = 4v_{\xi\xi}
\]
Après substitution, les termes croisés $v_{\xi\xi}$ s'annulent (cas parabolique), et on obtient la forme canonique :
\[
v_{\eta\eta} - v_\eta = 0
\]

\textbf{4.} En posant $w = v_\eta$, on a $w_\eta = w$, d'où $w = \Phi(\xi)\e^\eta$. Puis $v = \Phi(\xi)\e^\eta + \Psi(\xi)$, et :
\[
u(x,y) = \Phi(y-2x)\e^y + \Psi(y-2x)
\]
où $\Phi$ et $\Psi$ sont des fonctions arbitraires.
\end{boitecorrige}

\subsection*{Exercice 18 - Séries de Fourier \niveau{\moyen}}

\begin{boiteexercice}
Soit $f$ la fonction $2\pi$-périodique définie sur $[-\pi, \pi]$ par $f(x) = x$.
\begin{enumerate}
  \item Calculer les coefficients de Fourier $a_n$ et $b_n$.
  \item Écrire la série de Fourier de $f$.
  \item En utilisant le théorème de Parseval, calculer $\displaystyle\sum_{n=1}^{+\infty} \frac{1}{n^2}$.
  \item En déduire $\displaystyle\sum_{n=1}^{+\infty} \frac{(-1)^{n+1}}{n}$.
\end{enumerate}
\end{boiteexercice}

\begin{boitecorrige}
\textbf{1.} La fonction $f(x) = x$ est impaire sur $[-\pi,\pi]$, donc $a_n = 0$ pour tout $n$.
\[
b_n = \frac{1}{\pi}\int_{-\pi}^{\pi} x\sin(nx)\,\d x = \frac{2}{\pi}\int_0^{\pi} x\sin(nx)\,\d x
\]
Par intégration par parties : $b_n = \dfrac{2(-1)^{n+1}}{n}$.

\textbf{2.} Série de Fourier : $f(x) \sim \displaystyle 2\sum_{n=1}^{+\infty} \frac{(-1)^{n+1}}{n}\sin(nx)$.

\textbf{3.} Théorème de Parseval : $\dfrac{1}{\pi}\int_{-\pi}^{\pi}|f(x)|^2\,\d x = \dfrac{a_0^2}{2} + \sum_{n=1}^{+\infty}(a_n^2+b_n^2)$.
\[
\frac{1}{\pi}\int_{-\pi}^{\pi} x^2\,\d x = \frac{2\pi^2}{3} = \sum_{n=1}^{+\infty} \frac{4}{n^2}
\]
\[
\sum_{n=1}^{+\infty} \frac{1}{n^2} = \frac{\pi^2}{6}
\]

\textbf{4.} En $x = \pi/2$ (point de continuité) : $f(\pi/2) = \pi/2 = 2\sum_{n=1}^{+\infty}\frac{(-1)^{n+1}}{n}\sin(n\pi/2)$.

Cette série donne $\dfrac{\pi}{4} = 1 - \dfrac{1}{3} + \dfrac{1}{5} - \cdots = \sum_{n=0}^{+\infty}\dfrac{(-1)^n}{2n+1}$.
\end{boitecorrige}

\subsection*{Exercice 19 - Équation de Laplace en coordonnées polaires \niveau{\moyen}}

\begin{boiteexercice}
On cherche les solutions harmoniques (solutions de $\Delta u = 0$) dans le disque de rayon $R$ centré en $O$, en coordonnées polaires $(r,\theta)$.
\begin{enumerate}
  \item Écrire l'opérateur laplacien en polaires et chercher les solutions sous forme séparée $u(r,\theta) = f(r)g(\theta)$.
  \item Montrer que $g(\theta)$ est une combinaison de $\cos(n\theta)$ et $\sin(n\theta)$, et que $f(r) = r^n$ ou $r^{-n}$.
  \item Pour un problème intérieur régulier en $r=0$, donner la solution générale. Comment s'écrit-elle si $u(R,\theta) = h(\theta)$ (condition au bord de Dirichlet) ?
\end{enumerate}
\end{boiteexercice}

\begin{boitecorrige}

\textbf{1. Laplacien en polaires et séparation.}

En coordonnées polaires :
\[
\Delta u = \frac{\partial^2 u}{\partial r^2} + \frac{1}{r}\frac{\partial u}{\partial r} + \frac{1}{r^2}\frac{\partial^2 u}{\partial\theta^2}.
\]
En posant $u = f(r)g(\theta)$, l'équation $\Delta u = 0$ devient, après division par $fg$ :
\[
\frac{r^2 f'' + r f'}{f} = -\frac{g''}{g} = n^2 \quad (\text{constante de séparation}).
\]
On a posé la constante égale à $n^2$ pour assurer la $2\pi$-périodicité de $g$.

\textbf{2. Solutions angulaires et radiales.}

L'équation angulaire $g'' + n^2 g = 0$ donne $g(\theta) = A_n\cos(n\theta) + B_n\sin(n\theta)$ pour $n\ge 1$ et $g = C_0$ pour $n=0$.

L'équation radiale $r^2 f'' + r f' - n^2 f = 0$ est de type Euler. En cherchant $f = r^\alpha$, on obtient $\alpha(\alpha-1) + \alpha - n^2 = \alpha^2 - n^2 = 0$, soit $\alpha = \pm n$. Donc $f(r) = r^n$ ou $r^{-n}$ (et $f(r) = C + D\ln r$ pour $n=0$).

\textbf{3. Problème intérieur.}

Pour un disque contenant $r=0$, les solutions $r^{-n}$ et $\ln r$ divergent en $r=0$ : on les écarte. La solution générale harmonique dans le disque est :
\[
u(r,\theta) = a_0 + \sum_{n=1}^{+\infty} r^n\!\left(a_n\cos(n\theta) + b_n\sin(n\theta)\right).
\]
Avec la condition au bord $u(R,\theta) = h(\theta)$ (développable en série de Fourier $h(\theta) = \alpha_0 + \sum_n (\alpha_n\cos n\theta + \beta_n\sin n\theta)$), on identifie :
\[
a_0 = \alpha_0, \quad a_n = \frac{\alpha_n}{R^n}, \quad b_n = \frac{\beta_n}{R^n}.
\]
C'est la \emph{formule de Poisson} pour le disque : la solution est la projection harmonique du bord vers l'intérieur. Par exemple, $u(0,\theta) = a_0 = \frac{1}{2\pi}\int_0^{2\pi} h(\theta)\d\theta$ (valeur moyenne au bord), ce qui est la \emph{propriété de la moyenne} des fonctions harmoniques.

\reponse{$u(r,\theta) = a_0 + \sum_{n=1}^{+\infty} r^n(a_n\cos n\theta + b_n\sin n\theta)$}
\end{boitecorrige}

\subsection*{Exercice 20 - Équation des cordes vibrantes et séparation de variables \niveau{\difficile}}

\begin{boiteexercice}
On considère une corde vibrante de longueur $L$, fixée à ses extrémités, régie par l'équation d'onde :
\[
\frac{\partial^2 u}{\partial t^2} = c^2 \frac{\partial^2 u}{\partial x^2}, \quad u(0,t) = u(L,t) = 0, \quad u(x,0) = f(x), \quad \frac{\partial u}{\partial t}(x,0) = 0.
\]
\begin{enumerate}
  \item Chercher les solutions séparées $u(x,t) = X(x)T(t)$ et déterminer les modes propres.
  \item Écrire la solution générale satisfaisant les conditions aux limites.
  \item Donner la solution pour $f(x) = \epsilon\sin(\pi x/L)$ (mode fondamental).
\end{enumerate}
\end{boiteexercice}

\begin{boitecorrige}

\textbf{1. Séparation et modes propres.}

En posant $u = X(x)T(t)$ dans l'équation d'onde et en divisant par $XT$ :
\[
\frac{T''}{c^2 T} = \frac{X''}{X} = -\lambda \quad (\lambda > 0).
\]
L'équation spatiale $X'' + \lambda X = 0$ avec $X(0) = X(L) = 0$ donne les valeurs propres :
\[
\lambda_n = \frac{n^2\pi^2}{L^2}, \quad X_n(x) = \sin\!\left(\frac{n\pi x}{L}\right), \quad n=1,2,\ldots
\]
L'équation temporelle $T'' + c^2\lambda_n T = 0$ donne :
\[
T_n(t) = A_n\cos\!\left(\frac{n\pi c t}{L}\right) + B_n\sin\!\left(\frac{n\pi c t}{L}\right).
\]

\textbf{2. Solution générale.}

\[
u(x,t) = \sum_{n=1}^{+\infty}\!\left[A_n\cos\!\left(\frac{n\pi c t}{L}\right) + B_n\sin\!\left(\frac{n\pi c t}{L}\right)\right]\sin\!\left(\frac{n\pi x}{L}\right).
\]
Les coefficients sont déterminés par les conditions initiales :
\[
A_n = \frac{2}{L}\int_0^L f(x)\sin\!\left(\frac{n\pi x}{L}\right)\d x, \qquad B_n = \frac{2L}{n\pi c}\int_0^L g(x)\sin\!\left(\frac{n\pi x}{L}\right)\d x,
\]
où $g(x) = \partial_t u(x,0)$. Comme $g=0$ ici, $B_n = 0$.

\textbf{3. Mode fondamental $f(x) = \epsilon\sin(\pi x/L)$.}

On a $A_1 = \epsilon$ et $A_n = 0$ pour $n\ge 2$ (par orthogonalité des sinus). La solution se réduit au mode fondamental :
\[
u(x,t) = \epsilon\cos\!\left(\frac{\pi c t}{L}\right)\sin\!\left(\frac{\pi x}{L}\right).
\]
La corde vibre à la fréquence fondamentale $f_1 = c/(2L)$, sans aucune harmonique supérieure. C'est le principe de la corde de Melde : en excitant une corde à sa fréquence propre, on obtient un mode pur.

\reponse{$u(x,t) = \sum_n A_n\cos(n\pi c t/L)\sin(n\pi x/L)$ ; mode fondamental $f_1 = c/(2L)$}
\end{boitecorrige}

\subsection*{Exercice 21 - Fonctions de Green et solution d'une EDO linéaire \niveau{\difficile}}

\begin{boiteexercice}
On considère l'EDO avec conditions aux limites :
\[
u''(x) = f(x), \quad u(0) = u(1) = 0.
\]
\begin{enumerate}
  \item Déterminer la fonction de Green $G(x,\xi)$ solution de $G'' = \delta(x-\xi)$ avec $G(0,\xi) = G(1,\xi) = 0$.
  \item Vérifier que $u(x) = \int_0^1 G(x,\xi) f(\xi)\,\d\xi$ résout le problème.
  \item Application : $f(x) = -1$ (charge uniforme). Donner $u(x)$.
\end{enumerate}
\end{boiteexercice}

\begin{boitecorrige}

\textbf{1. Fonction de Green.}

Pour $x\neq\xi$, $G''=0$, donc $G$ est affine par morceaux. Avec $G(0,\xi)=0$ : $G(x,\xi) = A x$ pour $x<\xi$. Avec $G(1,\xi)=0$ : $G(x,\xi) = B(1-x)$ pour $x>\xi$.

Continuité en $x=\xi$ : $A\xi = B(1-\xi)$, soit $B = A\xi/(1-\xi)$.

Saut de la dérivée : en intégrant $G''=\delta$ sur un petit voisinage de $\xi$ :
\[
G'(\xi^+,\xi) - G'(\xi^-,\xi) = 1.
\]
Comme $G' = A$ pour $x<\xi$ et $G' = -B$ pour $x>\xi$ : $-B - A = 1$, soit $A = -1-B = -1-A\xi/(1-\xi)$, donc :
\[
A\!\left(1 + \frac{\xi}{1-\xi}\right) = -1 \quad\Rightarrow\quad A = -(1-\xi), \quad B = \xi.
\]
Finalement :
\[
G(x,\xi) = \begin{cases}
-\xi(1-x) & \text{si } x\ge\xi,\\
-x(1-\xi) & \text{si } x\le\xi.
\end{cases}
\]
On peut réécrire symétriquement : $G(x,\xi) = -\min(x,\xi)\,(1-\max(x,\xi))$. $G$ est symétrique en $x$ et $\xi$ (propriété d'auto-adjonction).

\textbf{2. Solution.}

L'opérateur $\partial_x^2$ appliqué à $u(x) = \int_0^1 G(x,\xi)f(\xi)\d\xi$ donne, par la propriété de la dérivée de $G$ :
\[
u''(x) = \int_0^1 G''(x,\xi) f(\xi)\d\xi = \int_0^1 \delta(x-\xi) f(\xi)\d\xi = f(x).
\]
Les conditions aux limites : $u(0) = \int_0^1 G(0,\xi) f(\xi)\d\xi = 0$ car $G(0,\xi)=0$, et de même $u(1)=0$.

\textbf{3. Application $f(x) = -1$.}

\[
u(x) = \int_0^1 G(x,\xi)(-1)\,\d\xi = (1-x)\int_0^x \xi\,\d\xi + x\int_x^1 (1-\xi)\,\d\xi.
\]
Calcul : $\int_0^x \xi\,\d\xi = x^2/2$ et $\int_x^1 (1-\xi)\,\d\xi = (1-x)^2/2$. Donc :
\[
u(x) = (1-x)\frac{x^2}{2} + x\frac{(1-x)^2}{2} = \frac{x(1-x)}{2}\,[x + (1-x)] = \frac{x(1-x)}{2}.
\]
C'est une parabole (élément fini P1 classique), maximale en $x=1/2$ avec $u_\text{max} = 1/8$.

\reponse{$G(x,\xi) = -\min(x,\xi)(1-\max(x,\xi))$ ; pour $f=-1$, $u(x) = x(1-x)/2$}
\end{boitecorrige}

\subsection*{Exercice 22 - Théorème de Stokes et champ à circulation conservative \niveau{\moyen}}

\begin{boiteexercice}
Soit le champ vectoriel $\vec{F}(x,y,z) = (yz,\, xz,\, xy)$ dans $\R^3$.
\begin{enumerate}
  \item Calculer $\rot\vec{F}$ et $\div\vec{F}$.
  \item Déduire que $\vec{F}$ dérive d'un potentiel scalaire $\phi$ que l'on déterminera.
  \item Vérifier le théorème de Stokes sur la portion de sphère de rayon $R$ au-dessus du plan $z=0$.
\end{enumerate}
\end{boiteexercice}

\begin{boitecorrige}

\textbf{1. Rotationnel et divergence.}

\[
\rot\vec{F} = \begin{vmatrix} \hat{e}_x & \hat{e}_y & \hat{e}_z \\
\partial_x & \partial_y & \partial_z \\
yz & xz & xy \end{vmatrix}
\]
\[
\rot\vec{F} = \bigl(\partial_y(xy) - \partial_z(xz),\;\partial_z(yz)-\partial_x(xy),\;\partial_x(xz)-\partial_y(yz)\bigr) = (x-x,\, y-y,\, z-z) = \vec{0}.
\]
Le champ est irrotationnel. Divergence : $\div\vec{F} = \partial_x(yz) + \partial_y(xz) + \partial_z(xy) = 0$.

\textbf{2. Potentiel scalaire.}

Comme $\rot\vec{F} = \vec{0}$ sur $\R^3$ (simplement connexe), il existe $\phi$ tel que $\vec{F} = \grad\phi$ :
\[
\partial_x\phi = yz \;\Rightarrow\; \phi = xyz + g(y,z).
\]
Alors $\partial_y\phi = xz + \partial_y g = xz \;\Rightarrow\; \partial_y g = 0$, donc $g = h(z)$. Enfin $\partial_z\phi = xy + h'(z) = xy \;\Rightarrow\; h'(z) = 0$, donc $h = \text{cste}$.

On choisit $\phi(x,y,z) = xyz$ (à constante additive près), et on vérifie : $\grad\phi = (yz, xz, xy) = \vec{F}$. $\checkmark$

\textbf{3. Théorème de Stokes.}

Le théorème affirme : $\oint_\mathcal{C}\vec{F}\cdot\d\vec{\ell} = \iint_\mathcal{S}(\rot\vec{F})\cdot\d\vec{S}$, où $\mathcal{C}$ est le bord de la surface $\mathcal{S}$ (ici, le cercle équatorial de rayon $R$ dans le plan $z=0$).

Comme $\rot\vec{F} = \vec{0}$, l'intégrale de surface est nulle, donc :
\[
\oint_\mathcal{C} \vec{F}\cdot\d\vec{\ell} = 0.
\]
Vérification directe : paramétrons $\mathcal{C}$ par $(R\cos\theta, R\sin\theta, 0)$, $\theta\in[0,2\pi]$. Sur $\mathcal{C}$, $\vec{F} = (0, 0, R^2\cos\theta\sin\theta)$. Le vecteur tangent $\d\vec{\ell} = (-R\sin\theta, R\cos\theta, 0)\,\d\theta$, donc $\vec{F}\cdot\d\vec{\ell} = 0$ (la composante $z$ de $\vec{F}$ ne contribue pas). L'intégrale est bien nulle, conformément au théorème de Stokes. C'est aussi une illustration directe du fait que $\vec{F} = \grad\phi$ implique $\vec{F}\cdot\d\vec{\ell} = \d\phi$, dont l'intégrale sur une courbe fermée est nulle (théorème fondamentale du gradient).

\reponse{$\phi(x,y,z) = xyz$, \quad $\oint_\mathcal{C}\vec{F}\cdot\d\vec{\ell} = 0$}
\end{boitecorrige}

\subsection*{Exercice 23 - Distribution de Dirac et dérivées faibles \niveau{\difficile}}

\begin{boiteexercice}
On rappelle qu'une distribution $T$ sur $\R$ est une forme linéaire continue sur l'espace des fonctions test $\mathcal{D}(\R)$ (fonctions $C^\infty$ à support compact). La distribution de Dirac est $\delta : \varphi\mapsto\varphi(0)$.
\begin{enumerate}
  \item Calculer la dérivée au sens des distributions de la fonction de Heaviside $H(x) = \mathbf{1}_{x>0}$.
  \item Montrer que $x\,\delta = 0$ au sens des distributions.
  \item Calculer $x\,\delta'$ et $\delta'$.
\end{enumerate}
\end{boiteexercice}

\begin{boitecorrige}

\textbf{1. Dérivée de Heaviside.}

La dérivée distributionnelle de $H$ est définie par $\langle H', \varphi\rangle = -\langle H, \varphi'\rangle$ pour toute fonction test $\varphi$ :
\[
\langle H', \varphi\rangle = -\int_0^{+\infty} \varphi'(x)\,\d x = -[\varphi(x)]_0^{+\infty} = -0 + \varphi(0) = \varphi(0) = \langle\delta, \varphi\rangle.
\]
Donc $H' = \delta$ au sens des distributions.

\textbf{2. Produit $x\,\delta$.}

Par définition, $\langle x\,\delta, \varphi\rangle = \langle\delta, x\,\varphi\rangle = (x\,\varphi)(0) = 0\cdot\varphi(0) = 0$. Donc $x\,\delta = 0$.

\textbf{3. Dérivée $\delta'$ et produit $x\,\delta'$.}

La dérivée de $\delta$ : $\langle\delta', \varphi\rangle = -\langle\delta, \varphi'\rangle = -\varphi'(0)$.

Pour le produit $x\,\delta'$, on utilise la formule de Leibniz distributions : $x\,\delta' = (x\,\delta)' - x'\,\delta = 0' - 1\cdot\delta = -\delta$. Vérification directe :
\[
\langle x\,\delta', \varphi\rangle = \langle\delta', x\,\varphi\rangle = -(x\,\varphi)'(0) = -\varphi(0) - 0\cdot\varphi'(0) = -\varphi(0) = \langle -\delta, \varphi\rangle.
\]
Donc $x\,\delta' = -\delta$.

Ces relations sont fondamentales en théorie des distributions. Elles justifient par exemple que $\delta(ax) = \delta(x)/|a|$ (changement d'échelle) et que la distribution de Dirac est l'élément neutre de la convolution : $f * \delta = f$.

\reponse{$H' = \delta$, \quad $x\,\delta = 0$, \quad $x\,\delta' = -\delta$}
\end{boitecorrige}

\subsection*{Exercice 24 - Équation de la chaleur sur un segment \niveau{\difficile}}

\begin{boiteexercice}
On considère l'équation de la chaleur sur le segment $[0,L]$ avec conditions de Dirichlet homogènes :
\[
\frac{\partial u}{\partial t} = D\,\frac{\partial^2 u}{\partial x^2}, \quad u(0,t) = u(L,t) = 0, \quad u(x,0) = f(x).
\]
\begin{enumerate}
  \item Chercher les solutions séparées $u = X(x)T(t)$ et identifier les modes propres.
  \item Donner la solution générale comme superposition de modes.
  \item Application : $f(x) = u_0\sin(\pi x/L)$. Donner $u(x,t)$ et commenter la décroissance.
\end{enumerate}
\end{boiteexercice}

\begin{boitecorrige}

\textbf{1. Séparation et modes.}

En posant $u = X(x)T(t)$ dans l'équation et en divisant par $DTX$ :
\[
\frac{T'}{DT} = \frac{X''}{X} = -\lambda, \quad \lambda > 0.
\]
Équation spatiale : $X'' + \lambda X = 0$ avec $X(0) = X(L) = 0$. Valeurs propres et fonctions propres :
\[
\lambda_n = \frac{n^2\pi^2}{L^2}, \quad X_n(x) = \sin\!\left(\frac{n\pi x}{L}\right), \quad n\ge 1.
\]
Équation temporelle : $T' + D\lambda_n T = 0$, donc $T_n(t) = \e^{-D\lambda_n t} = \e^{-n^2\pi^2 D t/L^2}$.

\textbf{2. Solution générale.}

\[
u(x,t) = \sum_{n=1}^{+\infty} a_n\,\e^{-n^2\pi^2 D t/L^2}\,\sin\!\left(\frac{n\pi x}{L}\right),
\]
avec $a_n = \dfrac{2}{L}\int_0^L f(x)\sin(n\pi x/L)\,\d x$ (coefficients de Fourier sinus de $f$).

\textbf{3. Application $f(x) = u_0\sin(\pi x/L)$.}

Par orthogonalité, seul $a_1 = u_0$ est non nul. La solution est :
\[
u(x,t) = u_0\,\e^{-\pi^2 D t/L^2}\,\sin\!\left(\frac{\pi x}{L}\right).
\]
Le profil spatial reste sinusoïdal (le mode fondamental), et l'amplitude décroît exponentiellement avec le temps caractéristique $\tau_1 = L^2/(\pi^2 D)$. Pour les modes supérieurs, $\tau_n = L^2/(n^2\pi^2 D) \ll \tau_1$ : les hautes fréquences spatiales sont dissipées beaucoup plus rapidement (le « lissage » de la chaleur). À temps long, seul le mode fondamental persiste.

\reponse{$u(x,t) = \sum_n a_n e^{-n^2\pi^2 Dt/L^2}\sin(n\pi x/L)$ ; temps caractéristique $\tau_n = L^2/(n^2\pi^2 D)$}
\end{boitecorrige}

\subsection*{Exercice 25 - Conformité et intégrale complexe : théorème de Jordan \niveau{\difficile}}

\begin{boiteexercice}
On veut calculer l'intégrale réelle :
\[
I = \int_{-\infty}^{+\infty} \frac{\d x}{(x^2+a^2)(x^2+b^2)}, \quad a, b > 0,\; a\neq b.
\]
\begin{enumerate}
  \item Identifier les pôles dans le demi-plan supérieur et justifier le choix du contour (demi-cercle supérieur).
  \item Appliquer le théorème des résidus pour calculer $I$.
  \item Vérifier par décomposition en éléments simples.
\end{enumerate}
\end{boiteexercice}

\begin{boitecorrige}

\textbf{1. Pôles et contour.}

L'intégrand $f(z) = 1/[(z^2+a^2)(z^2+b^2)]$ a quatre pôles simples : $z=\pm ia$, $z=\pm ib$. Dans le demi-plan supérieur : $z_1 = ia$ et $z_2 = ib$.

On ferme le contour par un demi-cercle supérieur $\mathcal{C}_R$ de rayon $R\to\infty$. Comme $|f(z)| \sim 1/R^4$ sur le demi-cercle et la longueur est $\pi R$, l'intégrale sur $\mathcal{C}_R$ tend vers 0 (lemme de Jordan). L'intégrale totale le long de l'axe réel + le demi-cercle vaut $2\pi i$ fois la somme des résidus.

\textbf{2. Résidus.}

Les résidus aux pôles simples $z_1 = ia$ et $z_2 = ib$ :
\[
\text{Res}(f, ia) = \frac{1}{(z+ia)(z^2+b^2)}\bigg|_{z=ia} = \frac{1}{(2ia)((ia)^2+b^2)} = \frac{1}{2ia(b^2-a^2)} = \frac{1}{2ia(b^2-a^2)}.
\]
De même : $\text{Res}(f, ib) = \dfrac{1}{2ib(a^2-b^2)} = -\dfrac{1}{2ib(b^2-a^2)}$.

\textbf{Calcul de $I$.}

\[
I = 2\pi i\!\left[\frac{1}{2ia(b^2-a^2)} - \frac{1}{2ib(b^2-a^2)}\right] = \frac{\pi}{b^2-a^2}\!\left[\frac{1}{a} - \frac{1}{b}\right] = \frac{\pi}{b^2-a^2}\cdot\frac{b-a}{ab} = \frac{\pi}{ab(a+b)}.
\]

\textbf{3. Vérification.}

Décomposition en éléments simples :
\[
\frac{1}{(x^2+a^2)(x^2+b^2)} = \frac{1}{b^2-a^2}\!\left[\frac{1}{x^2+a^2} - \frac{1}{x^2+b^2}\right].
\]
En intégrant de $-\infty$ à $+\infty$ et utilisant $\int_{-\infty}^{+\infty}\d x/(x^2+c^2) = \pi/c$ :
\[
I = \frac{1}{b^2-a^2}\!\left[\frac{\pi}{a} - \frac{\pi}{b}\right] = \frac{\pi(b-a)}{ab(b^2-a^2)} = \frac{\pi}{ab(a+b)}.
\]
On retrouve bien le résultat du calcul par résidus, ce qui valide la méthode.

\reponse{$I = \dfrac{\pi}{ab(a+b)}$}
\end{boitecorrige}

\ifdefined\inMainDoc\else\end{document}\fi
